\documentclass[10pt,a4paper]{amsart}
\usepackage[margin=19mm]{geometry}
\usepackage{amsmath,amssymb,mathtools}
\usepackage[scr=rsfs,cal=euler]{mathalfa}
\usepackage{xfrac}
\usepackage[unicode,hidelinks,bookmarks=false]{hyperref}
\newcommand{\ie}{i.e.\ }
\newcommand{\cf}{cf.\ }
\newcommand{\resp}{respectively\ }
\newcommand{\Subsection}{\subsection}
\providecommand{\nobreakdash}{\nobreak\hbox{-}}
\setlength{\emergencystretch}{2em}
\multlinegap0pt
\usepackage{bm}
\usepackage{bbm}
\usepackage{dsfont}
\usepackage{xspace}
\usepackage{cancel}
\usepackage{mathtools}
\usepackage{amsthm}
\makeatletter
\@ifundefined{thm}{\newtheorem{thm}{Thm}}{}
\@ifundefined{assump}{\newtheorem{assump}[thm]{Assumption}}{}
\@ifundefined{lemma}{\newtheorem{lemma}[thm]{Lemma}}{}
\makeatother
\DeclareMathOperator{\Gr}{Gr}
\newcommand{\GIDER}{\operatorname{GIDER}}
\newcommand{\GLERP}{\operatorname{GLERP}}
\newcommand{\bR}{\mathbb{R}}
\title[Geodesic Interpolation on the Grassmann Manifold: GLERP and ]{Geodesic Interpolation on the Grassmann Manifold: GLERP and Recursive GIDER Interpolants}
\author{Shingyu Leung}
\date{Source: arXiv:2606.16463v1 (CC BY 4.0)}
\begin{document}
\maketitle

\noindent\emph{Source and licence.} Geodesic Interpolation on the Grassmann Manifold: GLERP and Recursive GIDER Interpolants. Version arXiv:2606.16463v1 (CC BY 4.0), distributed under the Creative Commons Attribution 4.0 International licence, as recorded in the arXiv licence metadata for this exact version and shown on the abstract page https://arxiv.org/abs/2606.16463v1. Full rights notice: licenses/arxiv-2606.16463-CC-BY-4.0.txt.\par\smallskip

\noindent\textit{Editorial scope.} Counted unit: the Lemma whose source label is \texttt{lemma:local\_stability} in the article (source0.tex lines 540--556, source label \texttt{lemma:local\_stability}), delivered together with the article's own complete original proof of it (source0.tex lines 558--640, one \texttt{proof} environment of 83 lines). No mathematics is written, repaired or re-derived: statement and proof are the article's own text line for line. Required context retained uncounted: assump:normal (lines 532--534). Every definition, display and label of those lines is retained unchanged. Omitted from this package and not redistributed: 1--531, 535--539, 557--557, 641--1304. Nothing inside a retained excerpt is omitted or altered: every retained body is delivered exactly as the article's lines are written, so no omission receipt of any kind accompanies this package. Citations: the retained excerpts carry no citation command at all, so no bibliography entry of the article is transcribed and no reference is redistributed. Reference adaptation: none was needed and none was made; every reference inside the delivered unit resolves inside it, so no unresolved reference or citation is produced. Printed slips kept literal: no printed word, symbol, hypothesis or number of the counted statement or of its proof was altered. Layout and class: the article's own class is \texttt{article} with packages including grffile, latexsym, amssymb, enumerate, amsmath, epsfig, amsthm, geometry, setspace, color, graphicx, algorithm2e, empheq, bm; the wrapper is the lane's pinned \texttt{amsart} editorial wrapper at 10pt on A4, which restates the theorem-like environments the retained text needs and adds the article's own macro definitions that the retained text needs (\texttt{\textbackslash GIDER}, \texttt{\textbackslash GLERP}, \texttt{\textbackslash Gr}, \texttt{\textbackslash bR}), with extra wrapper packages (bm, bbm, dsfont, xspace, cancel, mathtools). The delivered numbering is the unit's own. Assets: no retained span contains a figure environment, a graphics-inclusion command, a PDF-page inclusion, a table or a TikZ picture; no image file is copied into this package, the package declares no asset dependency and no image file is redistributed with it. Licence: version arXiv:2606.16463v1 of this article is distributed under the Creative Commons Attribution 4.0 International licence, as recorded in the arXiv licence metadata for this exact version and shown on the abstract page https://arxiv.org/abs/2606.16463v1. A full rights notice is delivered at licenses/arxiv-2606.16463-CC-BY-4.0.txt. Characters: every retained excerpt is pure ASCII, so the delivered engine, XeLaTeX with its default Latin Modern text font, reports no missing character.
\begin{assump}[Local normal-neighborhood assumption]\label{assump:normal}
Let \(\gamma:[a,b]\to\Gr(r,m)\) be a smooth curve.  For the interpolation interval \([t_0,t_n]\), assume that \(h\) is sufficiently small so that the exact curve, the data points, and all intermediate \(\GLERP\) evaluations in the \(\GIDER_n\) recursion remain in a compact subset \(K\) of a common normal coordinate neighborhood \(U\subset\Gr(r,m)\).  In particular, all logarithm maps used by \(\GLERP\) are uniquely defined, and all geodesic interpolation and extrapolation steps in the recursion are well-defined.
\end{assump}

\begin{lemma}[Local coordinate stability]\label{lemma:local_stability}
Let Assumption \ref{assump:normal} hold, and let \(\phi:U\to\bR^d\), \(d=r(m-r)\), be a smooth coordinate chart on the normal neighborhood \(U\), with \(K\Subset U\).  Define
\[
    F_h(t)
    =
    \phi\left(\GIDER_n(\gamma(t_0),\ldots,\gamma(t_n);t)\right),
    \qquad t\in[t_0,t_n].
\]
If \(\gamma\in C^{n+1}\), then \(F_h\in C^{n+1}([t_0,t_n];\bR^d)\), and there exists a constant \(C\), independent of sufficiently small \(h\), such that
\[
    \max_{t\in[t_0,t_n]}
    \left\|
        \frac{d^{\,n+1}}{dt^{\,n+1}}F_h(t)
    \right\|
    \le C.
\]
\end{lemma}

\begin{proof}
The map \(\GLERP\) is a composition of the Grassmann logarithm map, scalar multiplication in a tangent space, and the Grassmann exponential map.  These maps are smooth on \(K\) because \(K\) lies inside a normal neighborhood and stays away from the cut locus.  Therefore every finite recursive composition appearing in \(\GIDER_n\) is smooth in \(t\).

It remains to justify that the derivatives do not blow up as \(h\to0\).  Write \(t=t_0+h\sigma\), where \(\sigma\in[0,n]\), and define the rescaled coordinate curve
\[
    \widetilde F_h(\sigma)
    =
    F_h(t_0+h\sigma).
\]
We shall show that
\[
    \frac{d^{\,q}}{d\sigma^{q}}\widetilde F_h(\sigma)
    =
    \mathcal{O}(h^q),
    \qquad q=1,\ldots,n+1,
\]
uniformly for \(\sigma\in[0,n]\).  This implies the desired estimate, because
\[
    \frac{d^{\,q}}{dt^q}F_h(t)
    =
    h^{-q}
    \frac{d^{\,q}}{d\sigma^q}\widetilde F_h(\sigma).
\]

The claim follows from the local expansion of geodesic interpolation.  Fix a point \(p\in K\), and use normal coordinates \(x=\phi(\mathcal{X})\) centered in a slightly larger coordinate neighborhood.  In these coordinates, the map
\[
    G(x,y;\lambda)
    =
    \phi\left(
        \GLERP(\phi^{-1}(x),\phi^{-1}(y);\lambda)
    \right)
\]
is smooth for \(x,y\in\phi(K)\) and for all values of \(\lambda\) generated by the fixed-order recursion.  Since \(G(x,x;\lambda)=x\) and the geodesic equation has smooth coefficients in normal coordinates, the small-displacement expansion has the form
\[
    G(x,y;\lambda)
    =
    (1-\lambda)x+\lambda y
    +
    \mathcal{R}(x,y;\lambda),
\]
where, for \(x\) and \(y\) sufficiently close,
\[
    \left\|\partial_\lambda^q \mathcal{R}(x,y;\lambda)\right\|
    \le
    C_q\|x-y\|^2,
    \qquad q=0,1,\ldots,n+1.
\]
The constants \(C_q\) are uniform because \(K\) is compact and the admissible \(\lambda\)-values are bounded for fixed \(n\).

The sampled data satisfy
\[
    \phi(\gamma(t_0+jh))
    =
    \phi(\gamma(t_0))
    +
    \sum_{\ell=1}^{n+1}
    \frac{(jh)^\ell}{\ell!}
    \frac{d^\ell}{dt^\ell}(\phi\circ\gamma)(t_0)
    +
    \mathcal{O}(h^{n+2}),
\]
uniformly for \(j=0,\ldots,n\).  In particular, neighboring sampled data differ by \(\mathcal{O}(h)\).  Applying the preceding expansion of \(G\) recursively, one obtains by induction over the recursion level \(k\) that each rescaled intermediate curve
\[
    \widetilde F_{j,k,h}(\sigma)
    =
    \phi\left(\mathcal{I}_j^k(t_0+h\sigma)\right)
\]
admits an expansion of the form
\[
    \widetilde F_{j,k,h}(\sigma)
    =
    \sum_{\ell=0}^{n+1}
    h^\ell P_{j,k,\ell}(\sigma)
    +
    \mathcal{O}(h^{n+2}),
\]
where the coefficients \(P_{j,k,\ell}\) are smooth functions on the fixed interval \([0,n]\), and, for \(\ell\le k\), they agree with the coefficients generated by the ordinary Neville interpolation of the Taylor coefficients of \(\phi\circ\gamma\).  In particular, the \(\sigma\)-derivatives of order \(q\) satisfy
\[
    \frac{d^q}{d\sigma^q}\widetilde F_{j,k,h}(\sigma)=\mathcal{O}(h^q),
    \qquad q=1,\ldots,n+1,
\]
uniformly for all \(j,k\) appearing in the recursion.  The induction uses that each normalized parameter \(\lambda_j^k(t_0+h\sigma)=(\sigma-j)/k\) is independent of \(h\), while the displacement between the two lower-level inputs to each \(\GLERP\) step is \(\mathcal{O}(h)\).  The remainder term \(\mathcal{R}\) is therefore at least quadratic in a quantity of size \(\mathcal{O}(h)\), and its derivatives with respect to \(\sigma\) satisfy the same scaling after the chain rule is applied.  Taking \(j=0\) and \(k=n\) gives the asserted bounds for \(\widetilde F_h\), and hence for \(F_h\).
\end{proof}

\end{document}
